\documentclass{blux}
\usepackage{mathptmx}
\usepackage[scaled=0.92]{helvet}
\usepackage{courier}

\usepackage{custom}

\newcommand{\dGAS}{$\delta$-GAS}
\newcommand{\Kinf}{\mathcal{K}_\infty}

\newcommand{\tr}{\to}
\newcommand{\trs}[1][\sigma]{\tr_{#1}}

\newcommand{\nqu}{n_{q,u}}
\newcommand{\nqx}{n_{q,x}}

\newcommand{\Ps}{\mathcal{P}}
\newcommand{\U}{\mathcal{U}}
\newcommand{\oV}{\opt{V}}
\newcommand{\sys}{S}
\newcommand{\As}{\mathcal{A}}
\newcommand{\as}{\alpha}
\newcommand{\Cs}{\mathcal{B}}
\newcommand{\cs}{\beta}

\title{Computing controlled invariant sets using semidefinite programming}
\author{Beno\^it Legat, Rapha\"el M. Jungers\\
        Universit\'e catholique de Louvain\\
        Avenue Georges Lemaitre, 4\\
        B-1348 Louvain-la-Neuve\\
        Belgium\\
    Email: {\tt benoit.legat@uclouvain.be},\\
           {\tt raphael.jungers@uclouvain.be}
        \and
    Paulo Tabuada\\
        Department of Electrical and Computer Engineering\\
        UCLA\\
    Email: {\tt tabuada@ee.ucla.edu}
}
\date{}

\begin{document}
\maketitle


The problem of computing a controlled invariant set is a paradigmatic challenge in the broad field of systems and control.
Indeed, it is for instance crucial in safety-critical applications,
such as the control of a platoon of vehicles or air traffic management; see for instance \cite{tomlin1998conflict},
where firm guarantees are needed on our ability to maintain the state in a safe region
(e.g., with a certain minimal distance between vehicles).
In other situations, the dynamical system might be too complicated to analyze exactly in every point of the state space,
but yet it can be possible to confine the state within a guaranteed set.


A set is \emph{controlled invariant} (sometimes also referred to as \emph{viable})
if, any trajectory whose initial point is in the set can be kept inside it by means of a proper control action.
Given a system with constraint specifications on the states and/or input, the controlled invariant set can be used
to determine initial states such that trajectories with these initial conditions are guaranteed to meet the specifications.
Moreover, in some situations, a state feedback control law can be derived from the knowledge of the controlled invariant set; see \cite{blanchini1999survey} for a survey.

The computation of invariant sets is usually achieved using either polyhedral computations or semidefinite programming.
If the system contains a control input, the computational complexity of the problem becomes even more challenging.
Indeed, this requires the computation of projections of polytopes (see e.g., the procedure p.~201 in \cite{blanchini2015set}) or the parametrization of a family of polytopes \cite{rakovic2010parameterized} when using polyhedral computations
and semidefinite programming techniques are not directly applicable.
The problem of polyhedral projection is well known to severely suffer from the curse of dimensionality. %



We give a method based on semidefinite programming to compute ellipsoidal invariant sets in \cite{legat2018computing}.
A key ingredient in our technique is that we work in the dual space of the geometric problem.
Given a discrete-time control system
\[ x_{k+1} = A x_k + B u_k, \quad x_k \in \mathcal{X} \]
we show that an ellipsoid
\[ \mathcal{E}_{P^{-1}} = \{\, x \in \R^n \mid x^\Tr P^{-1} x \leq 1 \,\} \]
is controlled invariant if and only if $\mathcal{E}_{P^{-1}} \subseteq \mathcal{X}$ and
\[ E A P A^\Tr E^\Tr \preceq E P E^\Tr \]
where $E$ is a projection on $\Image(B)^{\perp}$.

For a system with constraints on the input:
\[ x_{k+1} = A x_k + B u_k, \quad x_k \in \mathcal{X}, \quad u_k \in \mathcal{U} \]
we show that $\mathcal{E}_{P^{-1}}$ is controlled invariant if and only if
there exists a set
\[ \mathcal{E}_{Q^{-1}} = \{\, y \in \R^n \mid y^\Tr Q^{-1} y \leq 1 \,\} \]
such that $\mathcal{E}_{P^{-1}} \subseteq \mathcal{X}$, $\mathcal{E}_{Q^{-1}} \subseteq \mathcal{X} \times \mathcal{U}$ and
\[
    \begin{pmatrix}
        A & B
    \end{pmatrix}
    Q
    \begin{pmatrix}
        A & B
    \end{pmatrix}^\Tr
    \preceq P \preceq
    \begin{pmatrix}
        I & 0
    \end{pmatrix}
    Q
    \begin{pmatrix}
        I & 0
    \end{pmatrix}^\Tr.
\]


All these constraints are LMI-representable, and this allows us to introduce a new efficient procedure for the computation of controlled invariant sets.
Our general method is presented in \cite{legat2018computing} and can be applied to affine hybrid control system, we only present here the particular case of linear control systems.






\bibliographystyle{abbrv}
\bibliography{biblio}

\end{document}
